\documentclass{beamer}

\usepackage{pb-diagram}
%\usepackage{beamerthemesplit}
\usepackage[french]{babel}
\usepackage[utf8]{inputenc}
%\usetheme{default}
\useoutertheme{infolines}
\usepackage{graphicx, tikz}
\usepackage{comment}



\usepackage[utf8]{inputenc} % clavier et acccent
\usepackage[french]{babel} % langue francaise
\usepackage[french]{nomencl}
\usepackage[T1]{fontenc}
\makenomenclature
\usepackage{geometry} % package pour définir la zone du texte
\usepackage{fancyhdr} % package pour charger le style de la mise en page
\usepackage{amsfonts} % des symboles
\usepackage{pst-labo}% schémat labo 
\usepackage{amssymb}%  des symboles mathématiques
\usepackage{amsmath} % des symboles mathématiques
\usepackage{mathtools}
\usepackage{graphicx} % pour charger des figures
\usepackage[dvips]{epsfig}% pour charger des figures ".eps"
\usepackage{float}
\usepackage{multicol}  % pour écrire en plusieurs colonnes
\usepackage{pdfpages}
\usepackage{color}  % pour écrire en couleur
\usepackage{enumitem}
\usepackage{parskip}
\usepackage{amsthm}
\usepackage{hyperref} %pour insérer liens 
\usepackage{ stmaryrd } %pour les ensembles d'entiers
\usepackage{bm} %pour mettre formules en gras
\usepackage{nicematrix}

\usepackage{mathabx}

%\usepackage{mathabx}
\usepackage{dsfont}

\newcommand{\commPL}[1]{{\color{green}{#1}}}
\newcommand{\corrPL}[1]{{\color{blue}{#1}}}
\newcommand{\commCB}[1]{{\color{orange}{#1}}}
\newcommand{\commJB}[1]{{\color{red}{#1}}}

\setcounter{tocdepth}{2}

\graphicspath{{Figures/}}  % toutes les figures doivent être dans le dossier "Figures"
\usepackage{listings}

\usepackage{url}
\usepackage{array}
\usepackage{animate}
\usepackage{appendix}
\usepackage{bm}  % pour mettre des équations en gras
\usepackage[bottom]{footmisc} %JE L'AI RAJOUTE POUR METTRE DES NOTES DE BAS DE PAGE
\usepackage{hyperref} % pour des liens hypertexte

%\usepackage[citestyle=numeric]{biblatex}
%\addbibresource{bibliographie.bib}
%\AtBeginBibliography{\small}

% pour avoir la diagonale inverse dans les matrices
\newcommand{\revdots}{\mathinner
{\mkern1mu\raise1pt\vbox{\kern7pt\hbox{.}}
\mkern2mu\raise4pt\hbox{.}:
\mkern2mu\raise7pt\hbox{.}\mkern1mu}
}


\title{Formalisation d'un résultat sur les pavages en Rocq}

\author[L. BARROCHE]{Loryne BARROCHE}

%\institute[I2M]{Institut de mathématiques de Marseille - Universit\'e d'Aix-Marseille}

\institute[L2 MPCI]{L2 Mathématiques Physique Chimie Informatique - Universit\'e d'Aix-Marseille}

\date{Juin 2026}



%\newtheorem{thm}{Théorème}

%\newtheorem*{defn}{Définition}
%\newtheorem{prop}{Proposition}
%\newtheorem*{rk}{Remarque}
%\newtheorem*{ex}{Exemple}
%\newtheorem*{exs}{Exemples}
%\newtheorem*{mot}{Motivation}
%\newtheorem*{question}{Question}
%\newtheorem*{cor}{Corollaire}

\newcommand{\N}{\mathbb{N}}
\newcommand{\Z}{\mathbb{Z}}
\newcommand{\R}{\mathbb{R}}
\newcommand{\Q}{\mathbb{Q}}
\newcommand{\restrict}[2]{#1 \upharpoonright #2}
\newcommand{\arrows}[3]{\longrightarrow {#1}^{#2}_{#3}}
\newcommand{\Ur}{\mathbb{U}}
\newcommand{\titre}[1]{\begin{center} \textbf{#1} \end{center}}
\newcommand{\funct}[2]{#1 \longrightarrow #2}
\newcommand{\om}[1]{\textbf{#1},<^{\textbf{#1}}}
\newcommand{\m}[1]{\textbf{#1}}
\newcommand{\oc}[1]{\widetilde{\textbf{#1}},<^{\widetilde{\textbf{#1}}}}
\newcommand{\mc}[1]{\widetilde{\textbf{#1}}}
\newcommand{\iso}{\mathrm{iso}}
\newcommand{\s}{\mathbb{S}}
\newcommand{\dom}{\mathrm{dom}}
\newcommand{\Rt}{\overrightarrow{\mathcal{R}}}
\newcommand{\Cy}{S(2)}
\newcommand{\Aut}{\mathrm{Aut}}

\begin{document}

\begin{frame}
\titlepage

\end{frame}

\begin{frame}
\frametitle{L'utilisation de Rocq pour la formalisation de résultats}

%une petite présentation de ce qu'est Rocq

\end{frame}


\begin{frame}

%l'idée est de faire la démo voulue, et à des endroits où ça coince en Rocq donner des exemples de ce qu'il faut rajouter.

\begin{Definition}
Soit $C$ un ensemble de couleurs. On définit les tuiles de Wang en deux dimensions par des tuiles carrées le longueur 1 possédant une couleur sur chaque côté. Par la suite, nous considérerons chaque tuile comme un quadruplet $w$ = (Haut, Bas, Gauche, Droite) $\in C^4$.  On note $\tau$ un ensemble fini de tuiles de Wang en deux dimensions.
\end{Definition}

\begin{Definition}
Soient $c_{n | 1 \leq n \leq 8} \in C$. Soient $w_1 = (c_1, c_2, c_3, c_4)$ et $w_2 = (c_5, c_6, c_7, c_8)$. \\
On dit que $w_1$ est compatible à droite avec $w_2$ si et seulement si $c_4 = c_7$. On dit aussi que $w_2$ est compatible à gauche avec $w_1$.\\
On dit que $w_1$ est compatible en haut avec $w_2$ si et seulement si $c_1 = c_6$. On dit aussi que $w_2$ est compatible en bas avec $w_1$. \\
On dit que $w_1$ est compatible avec $w_2$ ssi $w_1$ est compatible à droite ou en haut avec $w_2$ ssi $w_2$ est compatible à gauche ou en bas avec $w_1$. On note $w_1 \text{ } \bm{\gamma_{\tau}} \text{ } w_2$.
\end{Definition}

\begin{Definition}
On dit que $P_{2D} : \mathds{Z}^2 \rightarrow \tau$ est un pavage de $\mathds{Z}^2$ ssi $\forall x \in \mathds{Z}^2, \exists w \in \tau, P_{2D}(x) = w$. Sans perte de généralité, on peut considérer que le centre de chaque tuile a des coordonnées relatives.
\end{Definition}

\begin{Definition}
Soit $P_{2D} : \mathds{Z}^2 \rightarrow \tau$ un pavage de $\mathds{Z}^2$ avec des tuiles de Wang. On dit que $P_{2D}$ est valide si et seulement si $\forall x = (x_1,x_2) \in \mathds{Z}^2, w_{(x_1, x_2)} \text{ } \bm{\gamma_{\tau}} \text{ } w_{(x_1+1, x_2)}, w_{(x_1, x_2)} \text{ } \bm{\gamma_{\tau}} \text{ } w_{(x_1-1, x_2)}, w_{(x_1, x_2)} \text{ } \bm{\gamma_{\tau}} \text{ } w_{(x_1, x_2+1)}, w_{(x_1, x_2)} \text{ } \bm{\gamma_{\tau}} \text{ } w_{(x_1, x_2-1)}$.
\end{Definition}

\begin{Definition}
On dit que $P_{2D} : \mathds{Z}^2 \rightarrow \tau$ est faiblement périodique ssi $\exists \Vec{u} = (x_u, y_u) \neq \Vec{0}$, $(x_u, y_u) \in \mathds{Z}^2$, $\forall x \in \mathds{Z}^2$, $\forall k \in \mathds{Z}$, {$P_{2D}(x+k\Vec{u}) = P_{2D}(x)$}\\
On dit que $\Vec{u}$ est un vecteur de périodicité.
\end{Definition}

\begin{Definition}
On dit que $P_{2D} : \mathds{Z}^2 \rightarrow \tau$ est fortement périodique ssi $\exists \Vec{u} = (x_u, y_u) \neq \Vec{0}$, $(x_u, y_u) \in \mathds{Z}^2, \exists \Vec{v} = (x_v, y_v) \neq \Vec{0}$, $(x_v, y_v) \in \mathds{Z}^2$, $\nexists a, b \in \mathds{R}$ tq $a\Vec{u} + b\Vec{v} = \Vec{0}$, et $\forall x \in \mathds{Z}^2$, $\forall k_1 \in \mathds{Z}$, $\forall k_2 \in \mathds{Z}$, $P_{2D}(x+k_1\Vec{u}+k_2\Vec{v}) = P_{2D}(x)$.
\end{Definition}

\NiceMatrixOptions{renew-dots,renew-matrix}

\begin{Definition}
Soit $P_{2D} : \mathds{Z}^2 \rightarrow \tau$. On forme un jeu de tuiles de Wang de longueur $l$ et de hauteur $h$ $T_{l, h}(\tau)$. Il est de la forme : 
$$
T_{l, h}(\tau) = 
\begin{bmatrix}
t_{max(0,h-1),0} & \Cdots & t_{max(0,h-1),j} & \Cdots & t_{max(0,h-1),max(0,l-1)}\\
\Vdots &  & \Vdots &  & \Vdots\\
t_{i, 0} & \Cdots & t_{i,j} & \Cdots & t_{i,max(0,l-1)}\\
\Vdots &  & \Vdots &  & \Vdots\\
t_{0,0} & \Cdots & t_{0,j} & \Cdots & t_{0,max(0,l-1)}
\end{bmatrix}
$$
On note le pavage d'une droite par un jeu de tuiles de Wang $T_{l, h}(\tau)$ avec l'application $P_{1D}^T : \mathds{Z} \rightarrow T_{l, h}(\tau)$
\end{Definition}

\begin{Definition}
Soit $P_{1D}^T : \mathds{Z} \rightarrow T_{l, h}(\tau)$. \\
On dit que deux jeux de tuiles $k_1$ et $k_2$ de $T_{l, h}(\tau)$ sont compatibles lorsque $\forall 0 \leq j \leq h-1$, $(P_{1D}^T(k_1))_{max(0,l-1), j} \text{ } \bm{\gamma_{\tau}} \text{ } (P_{1D}^T(k_2))_{0, j}$. On note $ P_{1D}^T(k_1) \text{ } \bm{\gamma_{T_{l, h}(\tau)}} \text{ } P_{1D}^T(k_2)$
\end{Definition}

\begin{Propriete}
Soit $P_{2D} : \mathds{Z}^2 \rightarrow \tau$ valide faiblement périodique de vecteur de périodicité $\Vec{u} = (x_u, y_u)$. Soit $P_{1D}^T : \mathds{Z} \rightarrow T_{|x_u|, |y_u|}(\tau)$. Alors $\forall 0 \leq i \leq max(0, |x_u| - 1)$, $\forall 0 \leq j \leq max(0, |y_u| - 1)$, $\forall k \in \mathds{Z}
$, $(P_{1D}^T(k))_{i,j} = P_{2D}(i + k|x_u|,j)$
\end{Propriete}

\begin{proof}
    $(P_{1D}^T(k))_{i,j}$ revient à considérer la tuile de Wang aux coordonnées $(i,j)$ dans le $k$ème jeu de tuiles. Dans $P_{2D}$, on avance de $k|x_u|$ à partir de l'origine pour atteindre ce jeu de tuiles. Puis on sélectionne la tuile de coordonnées $(i, j)$ dans ce jeu de tuiles.
\end{proof} %! pas ouf cette démo !

\begin{Lemme}
Soit $P_{2D} : \mathds{Z}^2 \rightarrow \tau$ valide faiblement périodique de vecteur de périodicité $\Vec{u} = (x_u, y_u)$. Alors $P_{1D}^T : \mathds{Z} \rightarrow T_{|x_u|, |y_u|}(\tau)$ est valide.
\end{Lemme}

\begin{proof}
    Soit $k \in \mathds{Z}$.
    Soit $ 0 \leq j \leq y_u$.
    D'après la propriété \ref{prop:p1dtp2d}, on a $(P_{1D}^T(k))_{|x_u|-1, j} = P_{2D}((k+1)|x_u|-1,j)$ et $(P_{1D}^T(k+1))_{0, j} = P_{2D}((k+1)|x_u|,j)$. Or, $P_{2D}$ est valide. Donc $P_{2D}((k+1)|x_u|-1,j) \text{ } \bm{\gamma_{\tau}} \text{ } P_{2D}((k+1)|x_u|,j)$. On en déduit que $(P_{1D}^T(k))_{|x_u|-1, j} \text{ } \bm{\gamma_{\tau}} \text{ } (P_{1D}^T(k+1))_{0, j}$
    Donc $P_{1D}^T(k) \text{ } \bm{\gamma_{T_{|x_u|, |y_u|}(\tau)}} \text{ } P_{1D}^T(k+1)$.
\end{proof}

\begin{Lemme}
$\forall P_{1D}^T : \mathds{Z} \rightarrow T_{l, h}(\tau)$, $\exists k_1 \in \mathds{Z}, \exists t \in \mathds{N}^*$, $P_{1D}^T(k_1) = P_{1D}^T(k_1+t)$. 
\end{Lemme}

\begin{proof}
    Soit $P_{1D}^T : \mathds{Z} \rightarrow T_{l, h}(\tau)$. $P_{1D}^T$ pave donc $\mathds{Z}$ avec des tuiles de $T_{l, h}(\tau)$. Comme $\tau$ est un ensemble fini de tuiles de Wang, $T_{l, h}(\tau)$ est aussi un ensemble fini. D'après le principe des tiroirs en version infinie, il existe au moins une tuile $w^T \in T_{l, h}(\tau)$ qui apparaît au moins deux fois dans $P_{1D}^T$. Donc il existe une position $k$ correspondant à la position de $w^T$, et un entier non nul $t$ tels que l'on retrouve $w^T$ $t$ tuiles plus loin, c'est-à-dire tels que l'on ait $P_{1D}^T(k) = P_{1D}^T(k+t)$.
\end{proof}

\begin{Lemme}
Soit $P_{1D}^T : \mathds{Z} \rightarrow T_{l, h}(\tau)$ tq $\exists k_1 \in \mathds{Z}, \exists t \in \mathds{N}^*$, $P_{1D}^T(k_1) = P_{1D}^T(k_1+t)$. Soit $D \in T_{tl,h}(\tau)$ le jeu de tuiles de dimensions $(t*l)*h$ dont la première tuile est $k_1$. On considère le pavage composé de l'unique tuile $P_{1D}^D : \mathds{Z} \rightarrow \{D\}$, défini par $\forall k \in \mathds{Z}$, $\forall 0 \leq i_1 \leq t*max(l - 1,0)$, $\forall 0 \leq j \leq max(h-1, 0)$, on décompose $i_1$ selon sa division entière par $l$ selon $i_1 = n*l+i$, $n, i \in \mathds{N}$ tels que $(P_{1D}^D(k))_{i_1,j} = (P_{1D}^T(n*k_1))_{i,j}$. Alors $P_{1D}^D$ est valide. On dit que $D$ est un domaine fondamental.


\end{Lemme}

\begin{proof}

Soit $P_{1D}^T : \mathds{Z} \rightarrow T_{l, h}(\tau)$. Soit $P_{1D}^D : \mathds{Z} \rightarrow \{D\}$. Soient $k \in \mathds{Z}$, $ 0 \leq i \leq max(l - 1,0)$, $ 0 \leq j \leq max(h-1, 0)$. Considérons les tuiles $k$ et $k+1$. Par définition, on a que $(P_{1D}^D(k))_{tl-1,j} = (P_{1D}^D(k))_{(t-1)l+(l-1),j} = (P_{1D}^T(k_1 + (t-1))_{l-1,j}$ et $(P_{1D}^D(k+1))_{0,j} = (P_{1D}^D(k))_{tl,j} = (P_{1D}^T(k_1+t))_{0,j}$. Or on sait que $P_{1D}^T$ est valide, donc par le \textbf{lemme \ref{lem:Tvalide}}, on a que $P_{1D}^T(k_1 + (t-1))_{l-1,j} \text{ } \bm{\gamma_{T_{l,h}(\tau)}} \text{ } (P_{1D}^T(k_1+t))_{0,j}$. On en déduit que $(P_{1D}^D(k)) \text{ } \bm{\gamma_{T_{tl,h}(\tau)}} \text{ } (P_{1D}^D(k+1))$. Donc $P_{1D}^D$ est valide.
\end{proof}

Remarque : le pavage $P_{1D}^D$ est périodique de période $t$.






%$(P_{1D}^D(k))_{i,j} = P_{2D}(k*l+i,j)$


\begin{Theoreme}
Soit $\tau$ un jeu fini de tuiles de Wang en deux dimensions.
$\exists P_{2D} : \mathds{Z}^2 \rightarrow \tau$ valide, $P_{2D}$ faiblement périodique ssi $\exists P_{2D}' : \mathds{Z}^2 \rightarrow \tau$ valide, $P_{2D}'$ fortement périodique.
\end{Theoreme}

\begin{proof}
    ($\Rightarrow$) Soit $P_{2D} : \mathds{Z}^2 \rightarrow \tau$ valide. Supposons que $P_{2D}$ faiblement périodique. D'après la \textbf{définition \ref{def:FaP}}, $\exists \Vec{u} = (x_u, y_u)$ un vecteur de périodicité. D'après le  \textbf{lemme \ref{lem:Tvalide}}, $P_{1D}^T : \mathds{Z} \rightarrow T_{|x_u|, |y_u|}(\tau)$ est valide. De plus, d'après le \textbf{lemme \ref{lem:Tperiode}}, $\exists k \in \mathds{Z} \exists t \in \mathds{N}^*, P_{1D}^T(k) = P_{1D}^T(k+t)$.\\
    - Si $y_u \neq 0$. !!!!!!!!!!!!!!!!! vérifier que tout ce qui est dit avant n'a pas besoin de cette hypothèse !!!!!!!!!!!!!!! Posons $\Vec{v} = (t|x_u|,0)$. Soit $D \in T_{t|x_u|, |y_u|}(\tau)$. Soit $P_{1D}^D : \mathds{Z} \rightarrow \{D\}$. Soient $0 \leq i \leq max(t|x_u| - 1,0)$ et $0 \leq j \leq max(|y_u|-1, 0)$. Soient $k_1, k_2 \in \mathds{Z}$. On pose $P_{2D}' : \mathds{Z}^2 \rightarrow \tau$ tq $P_{2D}'((i,j)+k_1\Vec{u}+k_2\Vec{v}) = (P_{1D}^D(x))_{i,j} \forall x \in \mathds{Z}$. Soit $x = (x_1, x_2) \in \mathds{Z}^2$. On décompose les coordonnées de $x$ selon $x_1 = k_{1x}(t|x_u|)+i$ et $x_2 = k_{2x}|y_u|+j$ avec $k_{1x} \in \mathds{Z}$, $k_{2x} \in \mathds{Z}$, $0 \leq i \leq max(t|x_u| - 1,0)$ et $0 \leq j \leq max(|y_u|-1, 0)$. On a alors $x=(i,j)+k_{1x}\Vec{u}+k_{2x}\Vec{v}$. On en déduit que $P_{2D}'(x+k_1\Vec{u}+k_2\Vec{v}) = P_{2D}'((i,j)+(k_{1x}+k_1)\Vec{u}+(k_{2x} + k_2)\Vec{v}) = P_{2D}'(i,j) = P_{2D}'(x)$. Donc $P_{2D}'$ est fortement périodique. De plus, il est valide selon $\Vec{u}$ par hypothèse, et selon $\Vec{v}$ d'après le \textbf{lemme \ref{lem:fondadomain}}. Donc $P_{2D}'(x)$ est valide.
    - Si $y_u = 0$.
    \newline
    ($\Leftarrow$) 
\end{proof}
    

\end{frame}

\end{document}